\subsection{An exact stopping criterion for the fixed action bank}
Let $Q=2^{80}$, and retain the saved grid of200 cells and106 actions.
The two arrays $H,h$ here are integer numerators of certified finite
lower defects, not exact values of the analytic double-sieve functions.
Each transition forms nonnegative weighted sums from the previous
generation, takes the same prescribed downward integer floors, and
retains the preceding entry by a maximum. Write this fixed map as $T$.
It is order-preserving and $T(X)\ge X$ componentwise.

If a finite execution from $X_0$ reaches $U$ with $T(U)=U$, then every
subsequent execution of this same fixed map equals $U$. More generally,
if $X_0\le U$ and $T(U)\le U$, monotonicity gives $T^j(X_0)\le U$ for
every $j$. Thus a finite exact stationary vector is an upper bound on
all further gains obtainable by this particular rounded recurrence.
It is not an upper bound on the true sieve functions, better quadrature,
another grid, or a larger action bank. No infinite analytic limit is
inserted into the prime-counting argument.

The first experiment retains $K=6.7618$ and starts from the adopted
312-generation array. After63 changing steps, one further step is
exactly stationary. The maximum remaining change in $H$ is below
$3.485\times10^{-16}$. The second experiment first replaces the action
baselines using the already proved $K'=6.761767$. For a saved certified
net baseline slope $a_j>0$, it adds only
\[
 \left\lfloor Q(6.7618-K')a_j\right\rfloor
\]
to the saved312 baseline numerator. These are the previously audited
net slopes: the reduction of a positive variable-kernel saving is
already subtracted. All weights, targets, physical ceilings, and actual
state restrictions remain unchanged. After118 changing steps, one
more step is exactly stationary, so only a finite430-generation table
is needed. These finite states inherit the established typed induction;
the check of stationarity does not substitute for that induction.

For a bound on the potential improvement in the already used $U_3,U_4$
occurrences, let $\delta_H$ be the maximum cell increment from the old
312 table to a specified final vector. At common endpoints the change
of the adjacent minimum is also at most $\delta_H$.
In both costs the residual logarithmic level satisfies $L\ge1/6$:
the first segment has $L=(1-t)/2$, the enhanced middle segment has
$L=57/64-5t/2$ up to $t=57/224$, and the last classical segment has
$L=1/2-t$ with $t\le1/3$.
Using $a=417/5000$ and $d=59/250$, the Chen-coefficient gain is at most
\[
 6\delta_H\left(\log\frac{1/3}{a}+\log\frac d a\right).
\]
The factor4 in each cost cancels the global Chen factor$1/4$.
Positive-series logarithmic upper bounds make this rationally checkable.
For the first arm it is less than $5.072\times10^{-15}$, providing a
precise reason not to continue iterating the unchanged bank merely to
gain more digits. The second arm's coarse bound is not a positive
gain certificate; its actual insertion must be computed separately.

For the second arm, the companion one-dimensional certificate pulls
each cell back under $s=\max(2,(A-Bt)/a)$ and integrates its constant
nonnegative difference exactly up to a positive logarithm lower bound:
\[
 \int_l^r\frac{\delta}{t(A-Bt)}\,dt
 =\frac{\delta}{A}\log\frac{r(A-Bl)}{l(A-Br)}.
\]
It retains the existing clamp at$s=2$ and the shared first segment's
multiplicity2. This is a replacement difference from the312 table;
the older248-to312 gain is not added again. All other table occurrences,
minimum comparisons and ordinary restoration terms stay at their
recorded choices. Any future global adoption must pay additional
finite approximation and compiler allowances and use an independently
verified finite recurrence and insertion certificate.
